\begin{trapbox}
	\begin{description}[style=nextline, leftmargin=2.8em, labelsep=0.5em, itemsep=0.35em]
		\item[\textbf{\textcolor{CTrap}{易错点一（混淆自变量）}}] 错误地将 $f(x+1)$ 的定义域理解为 $x+1$ 的范围。
		\item[\textbf{\textcolor{CTrap}{正确理解（区分自变量）：}}] 定义域是自变量 $x$ 的取值范围，与括号内表达式无关。
		\item[\textbf{\textcolor{CTrap}{错因分析（忽略变量代换）：}}] 没有抓住“定义域是自变量的取值集合”这一本质。
	\end{description}

	\medskip
	\begin{description}[style=nextline, leftmargin=2.8em, labelsep=0.5em, itemsep=0.35em]
		\item[\textbf{\textcolor{CTrap}{易错点二（对应法则误用）}}] 认为 $f(x)$ 和 $f(2x)$ 的定义域相同。
		\item[\textbf{\textcolor{CTrap}{正确理解（分开求定义域）：}}] 对应法则相同，但定义域需要根据各自自变量重新确定。
		\item[\textbf{\textcolor{CTrap}{错因分析（法则等同误解）：}}] 误以为 $f$ 相同则括号内整体范围相同。
	\end{description}

	\medskip
	\begin{description}[style=nextline, leftmargin=2.8em, labelsep=0.5em, itemsep=0.35em]
		\item[\textbf{\textcolor{CTrap}{易错点三（端点判断错误）}}] 解不等式时漏看端点是否包含，导致定义域区间开闭出错。
		\item[\textbf{\textcolor{CTrap}{正确理解（关注端点条件）：}}] 原定义域的端点是否包含，直接影响复合定义域端点的开闭。
		\item[\textbf{\textcolor{CTrap}{错因分析（疏忽端点开闭）：}}] 只关注数值，忽略了原区间的端点性质。
	\end{description}
\end{trapbox}
